% ============================================================
% nostarch-style.tex — Stile No Starch Press per libri tecnici
% ============================================================
% Caratteristiche: layout divertente ma professionale,
% font sans-serif principale, colori vivaci, box evidenti,
% codice con sfondo scuro opzionale.
% ============================================================

% === LINGUA ===
\usepackage[\BabelLang]{babel}
\usepackage[autostyle]{csquotes}

% === BIBLIOGRAFIA ===
\usepackage[
  backend=biber,
  style=numeric,
  sorting=nyt,
  urldate=long,
  maxbibnames=3,
  backref=true,
  backrefstyle=none
]{biblatex}

% === FONT (No Starch: New Century Schoolbook body, Futura heads) ===
\usepackage{fontspec}
\usepackage{amsmath}
\usepackage{amssymb}
\setmainfont{texgyreschola}[
  Extension=.otf,
  UprightFont=*-regular,
  BoldFont=*-bold,
  ItalicFont=*-italic,
  BoldItalicFont=*-bolditalic
]
\setsansfont{texgyreadventor}[
  Extension=.otf,
  Scale=0.90,
  UprightFont=*-regular,
  BoldFont=*-bold,
  ItalicFont=*-italic,
  BoldItalicFont=*-bolditalic
]
\setmonofont{InconsolataN}[
  Extension=.otf,
  Scale=0.85,
  UprightFont=*-Regular,
  BoldFont=*-Bold
]

% === LAYOUT PAGINA (7 × 9.25 poll. — formato No Starch) ===
\usepackage[
  paperwidth=7in,
  paperheight=9.25in,
  top=0.875in, bottom=0.875in,
  inner=1in, outer=0.875in,
  headheight=14pt,
  footskip=0.45in
]{geometry}

% === SPAZIATURA PARAGRAFI (No Starch: indent) ===
\setlength{\parindent}{1.5em}
\setlength{\parskip}{0pt plus 1pt}

% === COLORI (palette No Starch — nero/grigio con accenti vivaci) ===
\usepackage[dvipsnames,svgnames,table]{xcolor}
\definecolor{nostarch-black}{HTML}{1C1C1C}
\definecolor{nostarch-gray}{HTML}{4D4D4D}
\definecolor{nostarch-lightgray}{HTML}{8C8C8C}
\definecolor{nostarch-accent}{HTML}{D4451A}
\definecolor{code-bg}{HTML}{2B2B2B}
\definecolor{code-fg}{HTML}{E8E8E8}
\definecolor{code-light-bg}{HTML}{F0F0F0}
\definecolor{tip-green}{HTML}{2E994E}
\definecolor{warn-orange}{HTML}{E67A1A}
\definecolor{note-blue}{HTML}{3366CC}
\definecolor{link-blue}{HTML}{1A5DAA}
\definecolor{danger-red}{HTML}{CC3333}
\definecolor{gotcha-purple}{HTML}{7B2D8E}

% === GRAFICA E TABELLE ===
\usepackage{graphicx}
\usepackage{float}
\usepackage{booktabs}
\usepackage{longtable}
\usepackage{tabularx}
\usepackage{multirow}
\usepackage{array}
\usepackage{colortbl}

% === CODICE SORGENTE (No Starch: sfondo chiaro, bordo colorato) ===
\usepackage{listings}
\lstset{
  basicstyle=\footnotesize\ttfamily\color{nostarch-black},
  backgroundcolor=\color{code-light-bg},
  frame=l,
  framerule=2pt,
  rulecolor=\color{nostarch-accent},
  breaklines=true,
  breakatwhitespace=true,
  showstringspaces=false,
  tabsize=4,
  xleftmargin=1.2em,
  xrightmargin=0.5em,
  aboveskip=0.8em,
  belowskip=0.8em,
  numbers=left,
  numberstyle=\tiny\color{nostarch-lightgray},
  numbersep=8pt,
  keepspaces=true,
  keywordstyle=\color{nostarch-accent}\bfseries,
  stringstyle=\color{tip-green},
  commentstyle=\color{nostarch-lightgray}\itshape,
  columns=fullflexible,
  literate={€}{{\euro}}1 {→}{{$\rightarrow$}}1 {←}{{$\leftarrow$}}1
}
% Dark code variant
\lstdefinestyle{dark}{
  basicstyle=\footnotesize\ttfamily\color{code-fg},
  backgroundcolor=\color{code-bg},
  rulecolor=\color{nostarch-accent},
  keywordstyle=\color{yellow!80!white}\bfseries,
  stringstyle=\color{green!70!white},
  commentstyle=\color{gray!60!white}\itshape,
}
\lstdefinelanguage{yaml}{
  keywords={true,false,null,yes,no},
  sensitive=false,
  comment=[l]{\#},
  morestring=[b]",
  morestring=[b]',
}

% === INTESTAZIONI PIÈ DI PAGINA ===
\usepackage{fancyhdr}
\pagestyle{fancy}
\fancyhf{}
\fancyhead[LE]{\small\sffamily\color{nostarch-gray}\leftmark}
\fancyhead[RO]{\small\sffamily\color{nostarch-gray}\rightmark}
\fancyfoot[LE,RO]{\small\sffamily\bfseries\thepage}
\renewcommand{\headrulewidth}{1pt}
\renewcommand{\headrule}{\hbox to\headwidth{\color{nostarch-accent}\leaders\hrule height \headrulewidth\hfill}}

% === TITOLI CAPITOLO (No Starch: grande, grassetto, accento rosso) ===
\usepackage{titlesec}
\titleformat{\chapter}[display]
  {\normalfont\sffamily\bfseries}
  {\fontsize{80}{80}\selectfont\color{nostarch-accent!20}\thechapter}
  {-15pt}
  {\Huge\color{nostarch-black}}
  [\vspace{5pt}{\color{nostarch-accent}\titlerule[2.5pt]}]
\titlespacing*{\chapter}{0pt}{-30pt}{25pt}
\titleformat{\section}
  {\normalfont\Large\sffamily\bfseries\color{nostarch-black}}
  {\color{nostarch-accent}\thesection}{0.7em}{}
\titleformat{\subsection}
  {\normalfont\large\sffamily\bfseries\color{nostarch-gray}}
  {\color{nostarch-accent}\thesubsection}{0.7em}{}
\titleformat{\subsubsection}
  {\normalfont\normalsize\sffamily\bfseries\itshape\color{nostarch-gray}}
  {\thesubsubsection}{0.7em}{}

% === TITOLI PARTI ===
\titleformat{\part}[display]
  {\centering\normalfont\sffamily\bfseries}
  {\fontsize{60}{60}\selectfont\color{nostarch-accent!15}\partname\ \thepart}
  {10pt}
  {\Huge\color{nostarch-black}}
  [\vspace{4pt}{\color{nostarch-accent}\titlerule[3pt]}]

% === HYPERLINK ===
\usepackage{xurl}
\usepackage[
  colorlinks=true,
  linkcolor=nostarch-black,
  citecolor=nostarch-accent,
  urlcolor=link-blue,
  bookmarks=true,
  bookmarksnumbered=true
]{hyperref}

% === PROTEZIONE OVERFLOW ===
\tolerance=1000
\emergencystretch=1.5em
\setlength{\tabcolsep}{4pt}

% === SPAZIATURA LISTE ===
\usepackage{enumitem}
\setlist{nosep, leftmargin=1.8em, topsep=0.4em, itemsep=0.2em}

% === BOX ADMONITION ===
% ============================================================
% nostarch-boxes.tex — Box admonition stile No Starch Press
% ============================================================
% Stile: vivace, icone grandi, sfondo colorato, bordi evidenti.
% No Starch usa: Note, Warning, Gotcha!, Try It Yourself.
% ============================================================

\providecolor{tip-green}{HTML}{2E994E}
\providecolor{warn-orange}{HTML}{E67A1A}
\providecolor{note-blue}{HTML}{3366CC}
\providecolor{danger-red}{HTML}{CC3333}
\providecolor{gotcha-purple}{HTML}{7B2D8E}
\providecolor{nostarch-accent}{HTML}{D4451A}
\providecolor{nostarch-gray}{HTML}{4D4D4D}

\usepackage{fontawesome5}
\usepackage{tcolorbox}
\tcbuselibrary{skins,breakable}

% Note
\newtcolorbox{notebox}[1][]{
  colback=note-blue!6, colframe=note-blue,
  fonttitle=\sffamily\bfseries,
  title={\faInfoCircle\space \MakeUppercase{\LabelNote}},
  boxrule=1.5pt, arc=3pt, breakable,
  left=8pt, right=8pt,
  coltitle=white, colbacktitle=note-blue,
  toptitle=4pt, bottomtitle=4pt, #1
}

% Tip
\newtcolorbox{tipbox}[1][]{
  colback=tip-green!6, colframe=tip-green,
  fonttitle=\sffamily\bfseries,
  title={\faLightbulb\space \MakeUppercase{\LabelTip}},
  boxrule=1.5pt, arc=3pt, breakable,
  left=8pt, right=8pt,
  coltitle=white, colbacktitle=tip-green,
  toptitle=4pt, bottomtitle=4pt, #1
}

% Warning
\newtcolorbox{warnbox}[1][]{
  colback=warn-orange!6, colframe=warn-orange,
  fonttitle=\sffamily\bfseries,
  title={\faExclamationTriangle\space \MakeUppercase{\LabelWarning}},
  boxrule=1.5pt, arc=3pt, breakable,
  left=8pt, right=8pt,
  coltitle=white, colbacktitle=warn-orange,
  toptitle=4pt, bottomtitle=4pt, #1
}

% Danger / Gotcha!
\newtcolorbox{dangerbox}[1][]{
  colback=gotcha-purple!5, colframe=gotcha-purple,
  fonttitle=\sffamily\bfseries,
  title={\faBug\space \MakeUppercase{\LabelGotcha}},
  boxrule=1.5pt, arc=3pt, breakable,
  left=8pt, right=8pt,
  coltitle=white, colbacktitle=gotcha-purple,
  toptitle=4pt, bottomtitle=4pt, #1
}

% Checkpoint / Try It Yourself
\newtcolorbox{checkpointbox}[1][]{
  colback=nostarch-accent!5, colframe=nostarch-accent,
  fonttitle=\sffamily\bfseries,
  title={\faRocket\space \MakeUppercase{\LabelTryItYourself}},
  boxrule=1.5pt, arc=3pt, breakable,
  left=8pt, right=8pt,
  coltitle=white, colbacktitle=nostarch-accent,
  toptitle=4pt, bottomtitle=4pt, #1
}

% Version
\newtcolorbox{versionbox}[1][]{
  colback=purple!4, colframe=purple!50!black,
  fonttitle=\sffamily\bfseries,
  title={\faCogs\space \MakeUppercase{\LabelVersion}},
  boxrule=1pt, arc=3pt, breakable,
  left=8pt, right=8pt,
  coltitle=white, colbacktitle=purple!50!black,
  toptitle=4pt, bottomtitle=4pt, #1
}

% Cost
\newtcolorbox{costbox}[1][]{
  colback=Goldenrod!5, colframe=Goldenrod!80!black,
  fonttitle=\sffamily\bfseries,
  title={\faCoins\space \MakeUppercase{\LabelCost}},
  boxrule=1pt, arc=3pt, breakable,
  left=8pt, right=8pt,
  coltitle=white, colbacktitle=Goldenrod!80!black,
  toptitle=4pt, bottomtitle=4pt, #1
}

% Exercise
\newtcolorbox{praticabox}[1][]{
  colback=nostarch-gray!5, colframe=nostarch-gray,
  fonttitle=\sffamily\bfseries,
  title={\faPencilAlt\space \MakeUppercase{\LabelExercise}},
  boxrule=1.5pt, arc=3pt, breakable,
  left=8pt, right=8pt,
  coltitle=white, colbacktitle=nostarch-gray,
  toptitle=4pt, bottomtitle=4pt, #1
}


% === EURO ===
\newcommand{\euro}{€}

% === INDICE ANALITICO ===
\usepackage{makeidx}
\makeindex
